% ============================================================
%  UIDE – MCDA-A – Grupo 2 – Semana 2 – Reporte de ejercicio práctico
%  Formato: two-column IEEE-style con ajustes propios del grupo
% ============================================================
\documentclass[10pt,twocolumn]{article}

% ── Paquetes ────────────────────────────────────────────────
\usepackage[utf8]{inputenc}
\usepackage[T1]{fontenc}
\usepackage[spanish]{babel}
\usepackage{lmodern}
\usepackage{microtype}
\usepackage{geometry}
\usepackage{graphicx}
\usepackage{booktabs}
\usepackage{array}
\usepackage{amsmath}
\usepackage{amssymb}
\usepackage{hyperref}
\usepackage{xcolor}
\usepackage{fancyhdr}
\usepackage{titlesec}
\usepackage{enumitem}
\usepackage{caption}
\usepackage{float}
\usepackage{listings}
\usepackage{parskip}

% ── Geometría ───────────────────────────────────────────────
\geometry{
  a4paper,
  top=2cm, bottom=2.2cm,
  left=1.5cm, right=1.5cm,
  columnsep=0.6cm,
  headheight=14pt
}

% ── Colores ─────────────────────────────────────────────────
\definecolor{uideblue}{RGB}{0,70,127}
\definecolor{codered}{RGB}{180,0,0}
\definecolor{codegray}{RGB}{245,245,245}
\definecolor{codecomment}{RGB}{100,100,100}

% ── Hipervínculos ───────────────────────────────────────────
\hypersetup{
  colorlinks=true,
  linkcolor=uideblue,
  urlcolor=uideblue,
  citecolor=uideblue
}

% ── Títulos de secciones ────────────────────────────────────
\titleformat{\section}
  {\color{uideblue}\normalfont\bfseries\normalsize}
  {\thesection.}{0.4em}{}
\titleformat{\subsection}
  {\normalfont\bfseries\small}
  {\thesubsection.}{0.4em}{}
\titlespacing*{\section}{0pt}{6pt}{3pt}
\titlespacing*{\subsection}{0pt}{4pt}{2pt}

% ── Captions ────────────────────────────────────────────────
\captionsetup{font=small,labelfont=bf,skip=4pt}

% ── Código ──────────────────────────────────────────────────
\lstset{
  backgroundcolor=\color{codegray},
  basicstyle=\ttfamily\tiny,
  commentstyle=\color{codecomment},
  keywordstyle=\color{uideblue},
  language=Python,
  breaklines=true,
  frame=single,
  framerule=0.4pt,
  rulecolor=\color{gray!40},
  xleftmargin=4pt,
  xrightmargin=4pt,
  aboveskip=4pt,
  belowskip=4pt
}

% ── Encabezado / pie ────────────────────────────────────────
\pagestyle{fancy}
\fancyhf{}
\fancyhead[L]{\scriptsize\textit{MCDA-A – Grupo 2 | Aprendizaje profundo | Semana 2 | Reporte de ejercicio práctico}}
\fancyhead[R]{\includegraphics[height=20pt]{logo uide.png}}   % coloca tu logo aquí
\fancyfoot[C]{\thepage}
\renewcommand{\headrulewidth}{0.4pt}

% ── Sin sangría, espaciado entre párrafos ───────────────────
\setlength{\parindent}{0pt}
\setlength{\parskip}{4pt}

% ============================================================
\begin{document}

% ── Caja de metadatos (imitando el PDF de semana 1) ─────────
\twocolumn[{%
\begin{@twocolumnfalse}

  \noindent
  \begin{minipage}[t]{0.68\textwidth}
    \footnotesize Notebook del ejercicio Práctico
  \end{minipage}%
  \begin{minipage}[t]{0.32\textwidth}
    \raggedleft\footnotesize UIDE G2 W2 MCDA Notebook
  \end{minipage}

  \vspace{6pt}
  \begin{center}
    {\LARGE\bfseries Aplicación de Redes Neuronales Convolucionales (CNN) desde cero}\\[4pt]
    {\LARGE\bfseries para la Identificación de Patologías Pulmonares en Radiografías de Tórax}\\[10pt]
    {\normalsize Nataly Lanchimba / Dario Herrera / Jordy Cevallos}\\[2pt]
    {\small\textit{Universidad Internacional del Ecuador}}\\[2pt]
    {\small\textbf{Profesor:} Jefferson Rodriguez\quad|\quad\textbf{Semana 2}}
  \end{center}

  \vspace{6pt}
  \noindent
  \fbox{\begin{minipage}{\textwidth}
    \small\textbf{Código fuente}\\
    Proyecto en GitHub /
    Jupyter Notebook \href{https://github.com/TU-USUARIO/TU-REPO/blob/main/UIDE\_G2\_W2\_MCDA\_A\_NOV25\_Notebook.ipynb}%
    {\texttt{UIDE\_G2\_W2\_MCDA\_A\_NOV25\_Notebook.ipynb}}
    % ^^^ REEMPLAZA el URL anterior con el tuyo
  \end{minipage}}

  \vspace{10pt}
\end{@twocolumnfalse}
}]

% ============================================================
\section{Resumen del Problema}

El presente proyecto constituye la continuación del trabajo de la
Semana~1, en el que se utilizó un Perceptrón Multicapa (MLP) sobre el
dataset \textit{ChestX6: Multi-Class X-ray Dataset}.  En esta segunda
entrega se implementa una \textbf{Red Neuronal Convolucional (CNN)
desde cero}~--- sin pesos pre-entrenados~--- para el mismo problema de
clasificación automática de radiografías de tórax en seis categorías de
condiciones pulmonares/torácicas.

La motivación principal del cambio de arquitectura es que las imágenes
de rayos-X son datos espaciales 2D donde los patrones locales (nódulos,
infiltrados, cardiomegalia) son naturalmente detectables mediante filtros
convolucionales.  Las CNN son la solución óptima por cuatro razones:
\begin{itemize}[leftmargin=*,noitemsep,topsep=2pt]
  \item \textbf{Invarianza traslacional:} los filtros detectan patrones
        en cualquier zona de la imagen.
  \item \textbf{Jerarquía de características:} capas tempranas capturan
        bordes y texturas; capas profundas capturan estructuras
        anatómicas complejas.
  \item \textbf{Eficiencia paramétrica:} el peso compartido reduce
        drásticamente los parámetros frente a una red densa.
  \item \textbf{Evidencia empírica:} las CNN dominan los benchmarks de
        imagen médica (CheXNet, etc.).
\end{itemize}

Una RNN no sería adecuada (datos no secuenciales) ni un Transformer
puro (requeriría muchos más datos para converger desde cero).

% ============================================================
\section{Preparación del Dataset}

\subsection{Descripción del Dataset}

Se emplea el mismo \textit{ChestX6} de la Semana~1: 18\,036 imágenes de
radiografías de tórax en escala de grises, distribuidas en seis clases
de condiciones pulmonares: \textit{Normal}, \textit{Pneumonia-Bacterial},
\textit{Pneumonia-Viral}, \textit{COVID-19}, \textit{Tuberculosis} y
\textit{Emphysema}.

El dataset es relativamente equilibrado.  La clase con más imágenes es
\textit{Normal} (3\,271) y la con menos es \textit{Emphysema} (2\,550),
una diferencia de apenas el~28\,\% entre extremos, de modo que el
desbalance no es severo.  La división predefinida respeta una proporción
aproximada de 80/10/10, eliminando el riesgo de \textit{data leakage}.
Las características principales del dataset se resumen en el
Cuadro~\ref{tab:dataset}.

\begin{table}[H]
  \caption{Ficha técnica del dataset ChestX6}
  \label{tab:dataset}
  \centering\small
  \begin{tabular}{@{}lp{3.4cm}@{}}
    \toprule
    \textbf{Característica} & \textbf{Descripción} \\
    \midrule
    Número de muestras    & 18\,036 imágenes \\
    Modalidad             & Radiografía de tórax \\
    Espacio de color      & Escala de grises (1 canal nativo; RGB en CNN) \\
    Resolución nativa     & $224\times224$~px \\
    Resolución usada (CNN)& $128\times128$~px \\
    N.º de clases         & 6 \\
    Partición Train       & 14\,551 imágenes (80.7\,\%) \\
    Partición Validation  & 1\,748 imágenes (9.7\,\%) \\
    Partición Test        & 1\,737 imágenes (9.6\,\%) \\
    Tipo de tarea         & Clasificación multiclase \\
    \bottomrule
  \end{tabular}
\end{table}

\begin{figure}[H]
 \caption{Distribución de imágenes por clase del conjunto de datos}
    \centering
    \includegraphics[width=0.55\textwidth]{ficha_tecnica1.png}
    \label{fig:ficha_tecnica1}
\end{figure}

Como se puede ver en la Figura ~\ref{fig:ficha_tecnica1} presenta la distribución 
de imágenes por clase en los conjuntos de entrenamiento, validación y prueba. 
Se observa una distribución relativamente balanceada entre las diferentes categorías, 
con ligeras variaciones en el número de muestras disponibles. La clase Normal presenta 
la mayor cantidad de imágenes, mientras que Emphysema contiene el menor 
número de observaciones. Mientras que en validación y test se tienen proporciones
 similares en cada clase

\subsection{Preprocesamiento del Dataset}

Se aplicaron las siguientes transformaciones para adaptar las imágenes
al formato de entrada de la CNN:

\textbf{Reescalado:} de $224\times224$ a $128\times128$~píxeles.  Esta
reducción equilibra coste computacional y detalle anatómico suficiente
para la clasificación.

\textbf{Normalización:} división por 255, mapeando los valores de píxel
al rango $[0,1]$ para estabilizar el gradiente durante el
\textit{backpropagation} y acelerar la convergencia del optimizador.

\textbf{Aumento de datos} (solo en entrenamiento): \textit{flip}
horizontal, rotación leve ($\pm5$\,\%), zoom ($\pm10$\,\%), contraste
($\pm15$\,\%), brillo ($\pm15$\,\%) y traslación ($\pm10$\,\%).  Estas
transformaciones simulan variaciones naturales en la toma de
radiografías~--- diferencias de posicionamiento del paciente o ajustes
del equipo~--- sin distorsionar la información clínica relevante. Como se
puede ver a ocntinuación en la ~\ref{fig:preprocesamiento}:

\begin{figure}[H]
 \caption{Argumentación de la data}
    \centering
    \includegraphics[width=0.5\textwidth]{preprocesamiento.png}
    \label{fig:preprocesamiento}
\end{figure}

\textbf{Class weights:} se calcularon pesos de clase con
\texttt{compute\_class\_weight('balanced')} para compensar el leve
desbalance y evitar sesgo hacia las clases mayoritarias.

% ============================================================
\section{Implementación y Entrenamiento del Modelo}

\subsection{Arquitectura de la CNN}

La arquitectura está inspirada en VGG: bloques \textit{conv}$\times2$ +
\textit{MaxPooling}, sin pesos pre-entrenados.  El diseño es progresivo
(32~$\to$~64~$\to$~128~$\to$~256 filtros) para capturar características
de complejidad creciente:

\begin{enumerate}[leftmargin=*,noitemsep,topsep=2pt]
  \item \textbf{Bloque 1} ($32$ filtros): detecta bordes, gradientes y
        texturas de bajo nivel.
        \texttt{Conv2D(32)}$\times2$ + BN + ReLU + \texttt{MaxPool} +
        \texttt{Dropout(0.10)}\quad $128\to64$~px.
  \item \textbf{Bloque 2} ($64$ filtros): patrones intermedios,
        contornos anatómicos.
        \texttt{Conv2D(64)}$\times2$ + BN + ReLU + \texttt{MaxPool} +
        \texttt{Dropout(0.15)}\quad $64\to32$~px.
  \item \textbf{Bloque 3} ($128$ filtros): estructuras complejas,
        silueta cardíaca, nódulos, infiltrados.
        \texttt{Conv2D(128)}$\times2$ + BN + ReLU + \texttt{MaxPool} +
        \texttt{Dropout(0.20)}\quad $32\to16$~px.
  \item \textbf{Bloque 4} ($256$ filtros): representaciones de alto
        nivel específicas de cada clase.
        \texttt{Conv2D(256)}$\times2$ + BN + ReLU + \texttt{MaxPool} +
        \texttt{Dropout(0.25)}\quad $16\to8$~px.
  \item \textbf{Cabecera:} \texttt{GlobalAveragePooling2D} +
        \texttt{Dense(512)} + BN + ReLU + \texttt{Dropout(0.30)} +
        \texttt{Dense(256)} + BN + ReLU + \texttt{Dropout(0.35)} +
        \texttt{Dense(6, Softmax)}.
\end{enumerate}

El \textit{Global Average Pooling} sustituye a \texttt{Flatten} para
mejorar la generalización con menos parámetros.  La regularización
$L_2$ ($\lambda=10^{-4}$) se aplica en las capas densas.  El
Cuadro~\ref{tab:arch} resume la configuración completa.

\begin{table}[H]
  \caption{Configuración de la arquitectura CNN}
  \label{tab:arch}
  \centering\small
  \begin{tabular}{@{}llll@{}}
    \toprule
    \textbf{Bloque} & \textbf{Neuronas} & \textbf{Regularización} & \textbf{Salida} \\
    \midrule
    Conv Bloque 1  & $32\times2$   & BN + Dropout(0.10) & $64^2$  \\
    Conv Bloque 2  & $64\times2$   & BN + Dropout(0.15) & $32^2$  \\
    Conv Bloque 3  & $128\times2$  & BN + Dropout(0.20) & $16^2$  \\
    Conv Bloque 4  & $256\times2$  & BN + Dropout(0.25) & $8^2$   \\
    GAP            & ---           & ---                & 256     \\
    Dense 1        & 512           & BN+L2+Dropout(0.30)& 512     \\
    Dense 2        & 256           & BN+L2+Dropout(0.35)& 256     \\
    Salida         & 6             & Softmax            & 6       \\
    \midrule
    \multicolumn{2}{@{}l}{Función de pérdida} & \multicolumn{2}{l}{Sparse Cat. Crossentropy} \\
    \multicolumn{2}{@{}l}{Optimizador}        & \multicolumn{2}{l}{Adam + Cosine Decay} \\
    \multicolumn{2}{@{}l}{LR inicial}         & \multicolumn{2}{l}{$5\times10^{-4}$} \\
    \multicolumn{2}{@{}l}{Épocas máx.}        & \multicolumn{2}{l}{20} \\
    \multicolumn{2}{@{}l}{Batch size}         & \multicolumn{2}{l}{32} \\
    \bottomrule
  \end{tabular}
\end{table}

\subsection{Entrenamiento}

Se utilizó un \textit{scheduler} de tasa de aprendizaje
\textbf{Cosine Decay} (decaimiento desde $5\times10^{-4}$ hasta
$10^{-5}$) en lugar del \texttt{ReduceLROnPlateau} empleado en la
Semana~1, lo que produce una reducción suave y continua que evita
oscilaciones tardías.  Los \textit{callbacks} configurados fueron:

\begin{itemize}[leftmargin=*,noitemsep,topsep=2pt]
  \item \textbf{EarlyStopping} (paciencia=10, restaura mejores pesos,
        $\Delta_{\min}=10^{-4}$).
  \item \textbf{ModelCheckpoint} — guarda el mejor modelo según
        \texttt{val\_accuracy}.
  \item \textbf{CSVLogger} — registro del historial de entrenamiento.
\end{itemize}

Los pesos de clase fueron calculados con el objetivo de reducir el efecto del 
desbalance entre categorías durante el proceso de entrenamiento. Se observa 
que la clase Emphysema presenta el peso más alto (1.183), indicando una menor 
representación relativa dentro del conjunto de entrenamiento, mientras que la 
clase Normal presenta el menor peso (0.908), debido a su mayor número de 
muestras disponibles. Esta estrategia permite penalizar más los errores en 
clases menos representadas y contribuir a un aprendizaje más equilibrado 
del modelo.

\begin{table}[H]
\centering
\caption{Pesos de clase calculados para el entrenamiento del modelo}
\label{tab:class_weights}

\begin{tabular}{lc}
\hline
\textbf{Clase} & \textbf{Peso} \\
\hline
Covid-19 & 1.003 \\
Emphysema & 1.183 \\
Normal & 0.908 \\
Pneumonia-Bacterial & 1.010 \\
Pneumonia-Viral & 1.005 \\
Tuberculosis & 0.933 \\
\hline
\end{tabular}
\end{table}

% ============================================================
\section{Análisis del Preprocesamiento}

\subsection{Efecto del Data Augmentation}

Las transformaciones aplicadas al conjunto de entrenamiento
(flip horizontal, rotación leve, zoom, contraste y brillo) producen
imágenes visualmente diferentes pero clínicamente equivalentes: siguen
siendo radiografías válidas que un médico diagnosticaría de la misma
manera.  Este comportamiento es el deseado: aumentar artificialmente la
diversidad de entrenamiento sin introducir distorsiones que alteren el
diagnóstico.

\subsection{Diferencias morfológicas entre clases}

Entre las seis clases se observan diferencias morfológicas importantes:
\textit{Tuberculosis} presenta cavitaciones en lóbulos superiores y
fibrosis; \textit{COVID-19} y \textit{Emphysema} comparten opacidades
bilaterales difusas; \textit{Pneumonia-Bacterial} y
\textit{Pneumonia-Viral} presentan infiltrados intersticiales
bilaterales similares, lo que constituye el principal desafío de
clasificación del modelo.

% ============================================================
\section{Evaluación del Modelo}

El rendimiento del modelo se evaluó sobre el conjunto de prueba
utilizando \textit{accuracy}, precisión, \textit{recall} y
\textit{F1-score}, junto con la matriz de confusión y el análisis de
métricas por clase.  Adicionalmente se reporta la \textit{Top-2
Accuracy} como métrica complementaria.

\subsection{Métricas globales}

Las curvas de entrenamiento y validación muestran convergencia estable.
El análisis del \textit{gap} train-val confirma un ligero sobreajuste
(gap 2--10\,\%), aceptable para CNNs entrenadas desde cero.  El resumen
de métricas sobre el conjunto de test se presenta en el
Cuadro~\ref{tab:metrics}.

\begin{table}[H]
  \caption{Métricas de evaluación en el conjunto de test}
  \label{tab:metrics}
  \centering\small
  \begin{tabular}{@{}ll@{}}
    \toprule
    \textbf{Métrica} & \textbf{Valor} \\
    \midrule
    Test Loss              & (0.4259) \\
    Test Accuracy          & (0.8549) \\
    Top-2 Accuracy         & (0.9879) \\
    Precision (weighted)   & (0.8553) \\
    Recall (weighted)      & (0.8549) \\
    F1-Score (weighted)    & (0.8537) \\
    \bottomrule
  \end{tabular}
\end{table}

El análisis de las métricas globales presentado en la Tabla \ref{tab:metrics} evidencia un 
desempeño satisfactorio del modelo en el conjunto de prueba. La exactitud alcanzó un 
valor de 0.8549, indicando que aproximadamente el 85.49\% de las imágenes fueron 
clasificadas correctamente. Asimismo, el modelo obtuvo un F1-score ponderado de 
0.8537, mostrando un adecuado equilibrio entre precisión y sensibilidad. La precisión 
ponderada (0.8553) y el recall ponderado (0.8549) presentan valores similares, sugiriendo 
un comportamiento estable entre las diferentes categorías. Adicionalmente, la métrica T
op-2 Accuracy alcanzó un valor de 0.9879, lo que indica que en el 98.79\% de los casos 
la clase correcta se encuentra entre las dos predicciones con mayor probabilidad. 

\begin{table}[H]
\centering
\caption{Métricas de desempeño del modelo por clase}
\label{tab:metricas_clases}

\small
\begin{tabular}{p{3cm}cccc}
\hline
\textbf{Clase} & \textbf{Precision} & \textbf{Recall} & \textbf{F1-score}\\
\hline
Covid-19             & 0.8870 & 0.8633 & 0.8750 \\
Emphysema            & 0.8566 & 0.8840 & 0.8701 \\
Normal               & 0.9241 & 0.9733 & 0.9481 \\
Pneumonia-Bacterial  & 0.7717 & 0.6533 & 0.7076 \\
Pneumonia-Viral      & 0.6988 & 0.7733 & 0.7342 \\
Tuberculosis         & 1.0000 & 0.9930 & 0.9965 \\
\hline
\end{tabular}

\end{table}

\subsection{Análisis por clase}

Las métricas por clase permiten identificar las clases con mayor
dificultad de clasificación.  De manera consistente con los resultados
de la Semana~1, se espera que:

\begin{itemize}[leftmargin=*,noitemsep,topsep=2pt]
  \item \textbf{Tuberculosis} obtenga el F1-score más alto, dado que
        sus patrones radiológicos (cavitaciones en lóbulos superiores,
        fibrosis) son morfológicamente únicos.
  \item \textbf{Pneumonia-Viral} y \textbf{Pneumonia-Bacterial}
        registren el peor desempeño relativo, debido al solapamiento
        visual de infiltrados intersticiales bilaterales difusos.
\end{itemize}

\subsection{Matriz de confusión}

Las matrices de confusión (absoluta y normalizada) permiten identificar
los patrones de error del modelo.  Como se puede ver en la siguiente gráfica:

\begin{figure}[H]
 \caption{Matriz de confusión}
    \centering
    \includegraphics[width=0.45\textwidth]{confusion_matrix1.png}
    \label{fig:confusion_matrix1}
\end{figure}

Las principales confusiones esperadas se concentran entre:

\begin{itemize}[leftmargin=*,noitemsep,topsep=2pt]
  \item \textbf{Pneumonia-Bacterial} $\leftrightarrow$
        \textbf{Pneumonia-Viral}: solapamiento visual de infiltrados
        intersticiales.
  \item \textbf{COVID-19} $\leftrightarrow$ \textbf{Emphysema}: ambas
        presentan opacidades bilaterales difusas.

\end{itemize}

\subsection{Desempeño en entrenamiento y validación}

La curva de entrenamiento (azul) baja de forma suave y continua durante 
las 20 épocas, lo cual indica que el modelo aprendió de manera progresiva 
y estable. La curva de validación (rojo) es más irregular, con picos notables 
alrededor de las épocas 5--6 y 11, pero la tendencia general también es 
descendente. Que ambas converjan hacia valores similares ($\sim0.5$) al final 
es una buena señal: el modelo no está memorizando los datos de 
entrenamiento.

\begin{figure}[H]
    \centering
    \includegraphics[width=0.52\textwidth]{training_curves.png}
    \caption{Curvas de desempeño}
    \label{fig:training_curves1}
\end{figure}

La curva de entrenamiento sube de forma constante hasta alcanzar 0.9154 
al final. La de validación es más inestable al inicio (lo cual es normal con 
augmentation fuerte), pero se estabiliza a partir de la época 12 aproximadamente, 
terminando en 0.8358. El mejor epoch identificado fue el 18, con la mejor 
combinación entre el error del modelo y la exactitud de validación.

\subsection{Predicciones del modelo}

El análisis visual de las predicciones correctas e incorrectas evidencia que el modelo 
identifica adecuadamente patologías con patrones radiológicos distintivos, como 
tuberculosis y enfisema. Sin embargo, presenta dificultades para diferenciar 
enfermedades con características similares, particularmente entre neumonía viral y 
bacteriana. Además, se observaron casos críticos donde imágenes patológicas fueron 
clasificadas como normales con alta confianza, lo cual sugiere posibles problemas 
de generalización o desbalance entre clases y resalta la necesidad de mejorar el 
proceso de entrenamiento.

\begin{figure}[H]
    \centering
    \includegraphics[width=0.48\textwidth]{predicciones_ejemplos.png}
    \caption{Predicciones}
    \label{fig:predicciones_ejemplos}
\end{figure}

% ============================================================
\section{Conclusiones}

\begin{enumerate}[leftmargin=*,noitemsep,topsep=2pt]

\item \textbf{La CNN supera al MLP aprovechando la estructura espacial
de las imágenes.}  Mientras el MLP de la Semana~1 trataba cada imagen
como un vector plano de 4\,096 valores perdiendo toda la información
espacial (logrando un 86.13\,\% de \textit{accuracy}), la CNN extrae
jerarquías de características locales mediante filtros convolucionales,
lo que se traduce en una mejora directa sobre clases con patrones
morfológicos complejos.

\item \textbf{El diseño progresivo de filtros (32~$\to$~256) y el
Global Average Pooling son fundamentales.}  La arquitectura VGG-like
permite capturar desde bordes y texturas hasta estructuras anatómicas
complejas de manera jerárquica.  El \texttt{GlobalAveragePooling2D}
mejora la generalización frente al \texttt{Flatten} clásico, reduciendo
significativamente el número de parámetros en la cabecera clasificadora.

\item \textbf{El Cosine Decay produce convergencia más suave que el
ReduceLROnPlateau.}  Al decaer la tasa de aprendizaje de forma continua
y predecible, se evitan los saltos abruptos que pueden desestabilizar
el entrenamiento en épocas tardías.  Esto contribuye a curvas de
entrenamiento más estables.

\item \textbf{El model generaliza bien pero persisten dificultades en
las clases de neumonía.}  La confusión entre
\textit{Pneumonia-Bacterial} y \textit{Pneumonia-Viral} es médicamente
esperada: ambas condiciones comparten infiltrados intersticiales
bilaterales difusos, diferenciables principalmente por marcadores
clínicos y de laboratorio más que por apariencia radiológica.  Como
pasos futuros se recomienda implementar \textit{transfer learning} con
DenseNet-121 o EfficientNetB3, aplicar \textit{focal loss} para penalizar 
más los errores en estas clases difíciles, y explorar técnicas de aumento 
de datos específicas que exageren las diferencias entre patrones 
bacterianos y virales.

\item \textbf{Proyección hacia la Semana~3 (Transfer Learning):}
los modelos candidatos son DenseNet-121 (conexiones densas que capturan
gradientes finos de estructuras pulmonares), EfficientNetB3 (alto ratio
\textit{accuracy}/parámetros gracias al \textit{compound scaling}) y
ResNet-50 (\textit{baseline} robusto con \textit{skip connections} que
evitan el \textit{vanishing gradient}).  La estrategia consistirá en
\textit{fine-tuning} con pesos \texttt{imagenet} descongelando los
últimos bloques y añadiendo una nueva cabecera para 6 clases.

\end{enumerate}

% ============================================================
\section*{Referencias}

\begin{enumerate}[leftmargin=*,noitemsep,topsep=2pt,label={[\arabic*]}]
  \item TensorFlow Documentation (2024).
        \href{https://www.tensorflow.org/}{\texttt{https://www.tensorflow.org/}}
  \item Mohamed Asak (2024). \textit{ChestX6: Multi-Class X-ray Dataset}.
        Kaggle.
        \href{https://www.kaggle.com/datasets/mohamedasak/chest-x-ray-6-classes-dataset}%
        {\texttt{kaggle.com/datasets/mohamedasak/chest-x-ray-6-classes-dataset}}
  \item Rajpurkar et al.\ (2017). CheXNet: Radiologist-Level Pneumonia
        Detection on Chest X-Rays.\ \textit{arXiv:1711.05225}.
  \item Scikit-learn developers (2024). Neural network models (supervised).
        \href{https://scikit-learn.org/stable/modules/neural_networks_supervised.html}%
        {\texttt{scikit-learn.org}}
\end{enumerate}

\vfill
\noindent\rule{\columnwidth}{0.4pt}\\
{\scriptsize\textit{Aprendizaje profundo}\hfill
Universidad Internacional del Ecuador\hfill Semana 2}

\end{document}